%% ma: L10-B14-0070
%% lop: 10
%% chuong: 5
%% bai: 14
%% dang: 10.14.8
%% loai: TN4
%% muc: TH
%% dap_an: "D"
%% nguon: "chuyen-de-cac-so-dac-trung-cua-mau-so-lieu-khong-ghep-nhom-toan-10-chuong-trinh-moi.pdf (GV. Phan Nhật Linh) (P05), Bài 03 Dạng 2 Phần I Câu 14"
%% kien_thuc: "Lớp 10 Bài 13 (số trung bình), Bài 14 (phương sai, độ lệch chuẩn)"
%% trang_thai: chinh-thuc
%% ghi_chu: "Tự duyệt theo yêu cầu giáo viên 22/10/2026."
\begin{ex}
Bảng sau ghi lại điểm của $40$ học sinh trong bài kiểm tra môn Toán:
\begin{center}
\begin{tabular}{|c|c|c|c|c|c|c|c|c|c|}\hline
Điểm & $3$ & $4$ & $5$ & $6$ & $7$ & $8$ & $9$ & $10$ & Cộng\\\hline
Số học sinh & $2$ & $3$ & $7$ & $18$ & $3$ & $2$ & $4$ & $1$ & $40$\\\hline
\end{tabular}
\end{center}
Độ lệch chuẩn của mẫu số liệu gần bằng
\choice{$2{,}50$}{$2{,}33$}{$2{,}49$}{\True $1{,}58$}
\loigiai{$\overline x=6{,}1$; $s^2=\dfrac{1}{40}\sum n_i(x_i-6{,}1)^2=2{,}49$; $s=\sqrt{2{,}49}\approx1{,}58$. Chọn D.}
\end{ex}
